\chapter{空间格局与潜力释放条件分析}\label{ux7a7aux95f4ux683cux5c40ux4e0eux6f5cux529bux91caux653eux6761ux4ef6ux5206ux6790}

\section{空间分析方法}\label{ux7a7aux95f4ux5206ux6790ux65b9ux6cd5}

\subsection{核密度估计}\label{ux6838ux5bc6ux5ea6ux4f30ux8ba1}

在前文载体级指数计算的基础上，本章从空间尺度观察文化资源与旅游供给的错位。核密度估计（Kernel Density Estimation, KDE）分别用于计算 POI 与文化载体的空间热点。带宽统一设为 3.0 km，用于反映区县尺度下较为稳定的片区性集聚，而不是识别单个地块层面的微小差异。

\subsection{DBSCAN 聚类}\label{dbscan-ux805aux7c7b}

DBSCAN 聚类以 POI 为样本点，用于识别旅游供给的集聚边界。参数设置为ε=0.5 km，最小样本数为 30。该参数能够较清楚地区分桂城、大沥、狮山等城镇化片区中的 POI 高密度集聚，同时避免将零散点位误判为连续旅游片区。

\subsection{空间错位图叠加}\label{ux7a7aux95f4ux9519ux4f4dux56feux53e0ux52a0}

文化厚度与旅游热度在 500 m 网格尺度上进行归一化比较，并据此划分为四类空间状态：文化强而旅游弱的``沉睡潜力''、文化弱而旅游强的``空心景点''、文化与旅游均较强的``核心耦合''，以及两侧均较弱的``双低空白''。这一处理不只是画出热点，更重要的是判断文化资源和旅游产品在空间上是否真正发生对应。

\section{空间格局结果}\label{ux7a7aux95f4ux683cux5c40ux7ed3ux679c}

\subsection{POI 与文化载体的热点错位}\label{poi-ux4e0eux6587ux5316ux8f7dux4f53ux7684ux70edux70b9ux9519ux4f4d}

从核密度结果看，POI 高值区主要出现在桂城---大沥---狮山构成的东北部都市圈，并与广佛交界带、佛山国家高新区等现代城市功能空间较为重合。相较之下，文化载体高值区主要集中在西樵---九江---丹灶一带，与南海区传统农业、水利社会、古村落和非遗活动的历史地理基础更为接近。旅游供给和文化资源并非简单同向集聚，而是呈现出较稳定的方向性分离。

\begin{figure}[H]
\centering
\includegraphics[width=5.90556in,height=3.27018in]{ms_thesis/image12.png}
\caption{500 m 网格错位识别：基础口径与知识图谱口径对比}
\end{figure}

图 6.1 比较了 0 跳与 1 跳知识图谱口径下的网格错位结果。0 跳口径识别出 6 个核心耦合网格，1 跳口径识别出 8 个核心耦合网格。图中黄色星标用于标出核心耦合网格，其空间位置主要集中在西樵及少量东北部节点，说明文化厚度与旅游热度真正重叠的区域仍然有限。

\begin{figure}
\centering
\includegraphics[width=5.90556in,height=3.27018in]{ms_thesis/image13.png}
\caption{文化载体核密度图与 POI 核密度图对照}
\end{figure}

图 6.2 显示，文化载体核密度在西南片区形成连续高值，而 POI 核密度在桂城、大沥、狮山等片区更为突出。这个结果与第 5 章中 POI 镇街分布的统计结果相互印证，也为后文提出差异化提升策略提供空间依据。

\subsection{镇街对比}\label{ux9547ux8857ux5bf9ux6bd4}

\Needspace{12\baselineskip}
表 6.1　镇街尺度的 POI、核心文化载体与指数均值对比

{\def\LTcaptype{none} % do not increment counter
\begin{longtable}[]{@{}
  >{\centering\arraybackslash}p{(\linewidth - 12\tabcolsep) * \real{0.1429}}
  >{\centering\arraybackslash}p{(\linewidth - 12\tabcolsep) * \real{0.1429}}
  >{\centering\arraybackslash}p{(\linewidth - 12\tabcolsep) * \real{0.1429}}
  >{\centering\arraybackslash}p{(\linewidth - 12\tabcolsep) * \real{0.1429}}
  >{\centering\arraybackslash}p{(\linewidth - 12\tabcolsep) * \real{0.1429}}
  >{\centering\arraybackslash}p{(\linewidth - 12\tabcolsep) * \real{0.1429}}
  >{\centering\arraybackslash}p{(\linewidth - 12\tabcolsep) * \real{0.1429}}@{}}
\toprule\noalign{}
\endhead
\bottomrule\noalign{}
\endlastfoot
\textbf{镇街} & \textbf{POI 总量} & \textbf{核心文化载体数} & \textbf{CMI 均值} & \textbf{OAI 均值} & \textbf{THI 均值} & \textbf{MI 均值} \\
桂城街道 & 4,174 & 3 & 17.03 & 50.00 & 54.63 & 21.12 \\
里水镇 & 2,625 & 1 & 53.21 & 25.00 & 0.00 & −39.11 \\
狮山镇 & 2,233 & 3 & 22.62 & 41.67 & 15.41 & −16.74 \\
大沥镇 & 2,080 & 0 & --- & --- & --- & --- \\
丹灶镇 & 1,086 & 11 & 39.84 & 29.55 & 37.96 & 3.27 \\
西樵镇 & 1,055 & 26 & 52.71 & 27.88 & 34.95 & −5.35 \\
九江镇 & 259 & 15 & 36.71 & 35.00 & 35.85 & 0.00 \\
\end{longtable}
}

表 6.1 显示，桂城街道的 POI 总量最高，MI 均值为正，呈现旅游供给较强而文化记忆和载体支撑相对不足的特征。西樵镇核心文化载体数最多，CMI 均值也较高，但 THI 均值低于桂城和大沥，表现为``文化较厚、旅游转化不足''。九江镇核心文化载体数不少，但 POI 数量仅 259 条，旅游供给规模偏小。大沥镇在核心指数样本中缺少可用于计算 CMI、OAI 和 THI 的明确文化载体，因此本表不对其指数均值作判断，后文通过官方资源扩展层再作补充诊断。

\begin{figure}
\centering
\includegraphics[width=5.90556in,height=3.3311in]{ms_thesis/image14.png}
\caption{镇街尺度文化厚度、旅游热度与错位类型对比}
\end{figure}

\subsection{知识图谱 0 跳 / 1 跳口径复核}\label{ux77e5ux8bc6ux56feux8c31-0-ux8df3-1-ux8df3ux53e3ux5f84ux590dux6838}

为避免空间判断仅停留在 POI 与载体坐标叠加层面，本文在 500 m 网格尺度上引入知识图谱路径复核。0 跳口径表示文化载体直接命中的图谱实体集合；1 跳口径则在此基础上纳入与这些实体直接相邻的人物、事件、地点、文献和技艺等信息。旅游热度仍由 POI 密度、平均评分与评论量构成，来源独立于文化侧；错位值仍表示旅游热度相对于文化厚度的差异。

\Needspace{12\baselineskip}
表 6.2　500 m 网格 0 跳 / 1 跳错位分类对比

{\def\LTcaptype{none} % do not increment counter
\begin{longtable}[]{@{}
  >{\centering\arraybackslash}p{(\linewidth - 10\tabcolsep) * \real{0.1250}}
  >{\centering\arraybackslash}p{(\linewidth - 10\tabcolsep) * \real{0.1750}}
  >{\centering\arraybackslash}p{(\linewidth - 10\tabcolsep) * \real{0.1750}}
  >{\centering\arraybackslash}p{(\linewidth - 10\tabcolsep) * \real{0.1750}}
  >{\centering\arraybackslash}p{(\linewidth - 10\tabcolsep) * \real{0.1750}}
  >{\centering\arraybackslash}p{(\linewidth - 10\tabcolsep) * \real{0.1750}}@{}}
\toprule\noalign{}
\endhead
\bottomrule\noalign{}
\endlastfoot
\textbf{口径} & \textbf{双低空白} & \textbf{一般地带} & \textbf{沉睡潜力} & \textbf{空心景点} & \textbf{核心耦合} \\
0 跳 & 3,870 & 651 & 99 & 36 & 6 \\
1 跳 & 3,841 & 609 & 170 & 34 & 8 \\
\end{longtable}
}

南海区有效网格共 4,662 个。表 6.2 显示，知识图谱从 0 跳扩展到 1 跳后，沉睡潜力网格由 99 个增加到 170 个，核心耦合网格由 6 个增加到 8 个，空心景点网格则略降为 34 个。这个变化说明，部分地点虽然没有直接命中文化载体名称，但可以通过图谱关系连接到相关人物、事件、技艺和文献，从而显化出原先被低估的文化厚度。

\Needspace{12\baselineskip}
表 6.3　镇街尺度 1 跳知识图谱口径摘要

{\def\LTcaptype{none} % do not increment counter
\small
\setlength{\tabcolsep}{2.2pt}
\renewcommand{\arraystretch}{1.12}
\begin{longtable}[]{@{}
  >{\centering\arraybackslash}p{0.130\linewidth}
  >{\centering\arraybackslash}p{0.075\linewidth}
  >{\centering\arraybackslash}p{0.090\linewidth}
  >{\centering\arraybackslash}p{0.075\linewidth}
  >{\centering\arraybackslash}p{0.100\linewidth}
  >{\centering\arraybackslash}p{0.070\linewidth}
  >{\centering\arraybackslash}p{0.100\linewidth}
  >{\centering\arraybackslash}p{0.080\linewidth}
  >{\centering\arraybackslash}p{0.080\linewidth}
  >{\centering\arraybackslash}p{0.080\linewidth}@{}}
\toprule\noalign{}
\endhead
\bottomrule\noalign{}
\endlastfoot
\textbf{镇街} & \textbf{网格数} & \textbf{文化载体数} & \textbf{POI 数} & \textbf{C（1 跳）均值} & \textbf{T 均值} & \textbf{M（1 跳）均值} & \textbf{1 跳沉睡} & \textbf{1 跳空心} & \textbf{1 跳核心} \\
西樵镇 & 715 & 95 & 1,310 & 11.78 & 6.55 & −5.23 & 89 & 4 & 6 \\
九江镇 & 386 & 34 & 531 & 6.17 & 5.70 & −0.47 & 23 & 0 & 1 \\
丹灶镇 & 580 & 32 & 715 & 4.36 & 5.07 & 0.71 & 35 & 0 & 0 \\
桂城街道 & 333 & 29 & 3,859 & 2.15 & 22.28 & 20.13 & 7 & 16 & 0 \\
大沥镇 & 397 & 14 & 2,593 & 1.73 & 17.18 & 15.45 & 7 & 6 & 1 \\
狮山镇 & 1,403 & 3 & 2,567 & 0.32 & 5.72 & 5.40 & 6 & 4 & 0 \\
里水镇 & 619 & 4 & 1,449 & 0.37 & 7.59 & 7.22 & 3 & 3 & 0 \\
\end{longtable}
}

表 6.3 显示，西樵---九江---丹灶仍是 1 跳沉睡潜力的主体，三镇合计 147 个沉睡潜力网格，占全区 1 跳沉睡潜力网格的 86.5\%。其中，西樵镇同时拥有 6 个核心耦合网格，既是文化潜力集中区，也是少数已经形成文化与旅游高值重叠的片区。桂城街道虽然旅游热度最高，但 1 跳口径下仍有 16 个空心景点网格，其问题不是``没有旅游''，而是现代旅游供给与本土文化叙事之间仍缺少稳定连接。

\section{潜力释放条件的相关性分析}\label{ux6f5cux529bux91caux653eux6761ux4ef6ux7684ux76f8ux5173ux6027ux5206ux6790}

在前文错位分层的基础上，本节观察 CMI、OAI 与 THI 三类变量之间的关系。这里不试图证明某一变量对另一变量具有因果影响，而是判断文化记忆、官方认证和旅游热度是否在同一空间载体上形成一致支撑。若三者同步较高，该载体具备较强优先深化条件；若某一侧明显偏弱，后续干预则应围绕短板展开。

\subsection{载体级相关矩阵（n = 165）}\label{ux8f7dux4f53ux7ea7ux76f8ux5173ux77e9ux9635n-165}

\Needspace{14\baselineskip}
表 6.4　载体级指标相关矩阵（Pearson, n = 165）

{\def\LTcaptype{none} % do not increment counter
\begin{longtable}[]{@{}
  >{\centering\arraybackslash}p{(\linewidth - 14\tabcolsep) * \real{0.1250}}
  >{\centering\arraybackslash}p{(\linewidth - 14\tabcolsep) * \real{0.1250}}
  >{\centering\arraybackslash}p{(\linewidth - 14\tabcolsep) * \real{0.1250}}
  >{\centering\arraybackslash}p{(\linewidth - 14\tabcolsep) * \real{0.1250}}
  >{\centering\arraybackslash}p{(\linewidth - 14\tabcolsep) * \real{0.1250}}
  >{\centering\arraybackslash}p{(\linewidth - 14\tabcolsep) * \real{0.1250}}
  >{\centering\arraybackslash}p{(\linewidth - 14\tabcolsep) * \real{0.1250}}
  >{\centering\arraybackslash}p{(\linewidth - 14\tabcolsep) * \real{0.1250}}@{}}
\toprule\noalign{}
\endhead
\bottomrule\noalign{}
\endlastfoot
\textbf{指标} & \textbf{CMI} & \textbf{OAI} & \textbf{THI} & \textbf{MI} & \textbf{POI（500 m）} & \textbf{评分（500 m）} & \textbf{评论（500 m）} \\
CMI & 1.000 & 0.173 & 0.086 & −0.495 & −0.031 & 0.048 & 0.114 \\
OAI & 0.173 & 1.000 & 0.165 & −0.440 & 0.069 & 0.145 & 0.041 \\
THI & 0.086 & 0.165 & 1.000 & 0.682 & 0.693 & 0.859 & 0.506 \\
MI & −0.495 & −0.440 & 0.682 & 1.000 & 0.538 & 0.597 & 0.332 \\
POI（500 m） & −0.031 & 0.069 & 0.693 & 0.538 & 1.000 & 0.423 & 0.360 \\
评分（500 m） & 0.048 & 0.145 & 0.859 & 0.597 & 0.423 & 1.000 & 0.266 \\
评论（500 m） & 0.114 & 0.041 & 0.506 & 0.332 & 0.360 & 0.266 & 1.000 \\
\end{longtable}
}

\begin{figure}[p]
\centering
\includegraphics[width=0.88\textwidth]{ms_thesis/image15.png}
\caption{载体级相关矩阵热图（7 变量 × 7 变量）}
\end{figure}

载体级相关矩阵给出三点观察。第一，CMI 与 MI 呈中等负相关（r = −0.495），OAI 与 MI 也呈中等负相关（r = −0.440），文化记忆越深、官方认证越高的载体，越容易表现为旅游热度相对不足。第二，500 m 范围内 POI 数量与 CMI、OAI 的相关系数分别为 −0.031 和 0.069，几乎不相关，说明 POI 的空间密度并没有自然跟随文化家底分布。第三，THI 与评分、POI 数量和评论量的相关性较高，这与其指标构成有关，可以说明指标内部逻辑基本符合预期，但不能被理解为外部独立验证。

从潜力解释看，CMI、OAI 和 THI 三者并不是简单相互替代的关系。CMI 较高但 THI 偏低，提示文化资源尚未转化为稳定旅游热度；THI 较高但 CMI、OAI 偏低，提示已有客流和场地基础，但地方文化叙事仍需补强；三者同时较高的载体，则应理解为具备优先深化潜力的对象。

\subsection{镇街级相关矩阵（n = 7）}\label{ux9547ux8857ux7ea7ux76f8ux5173ux77e9ux9635n-7}

\Needspace{12\baselineskip}
表 6.5　镇街级指标相关矩阵（Pearson, n = 7）

{\def\LTcaptype{none} % do not increment counter
\begin{longtable}[]{@{}
  >{\centering\arraybackslash}p{(\linewidth - 14\tabcolsep) * \real{0.1250}}
  >{\centering\arraybackslash}p{(\linewidth - 14\tabcolsep) * \real{0.1250}}
  >{\centering\arraybackslash}p{(\linewidth - 14\tabcolsep) * \real{0.1250}}
  >{\centering\arraybackslash}p{(\linewidth - 14\tabcolsep) * \real{0.1250}}
  >{\centering\arraybackslash}p{(\linewidth - 14\tabcolsep) * \real{0.1250}}
  >{\centering\arraybackslash}p{(\linewidth - 14\tabcolsep) * \real{0.1250}}
  >{\centering\arraybackslash}p{(\linewidth - 14\tabcolsep) * \real{0.1250}}
  >{\centering\arraybackslash}p{(\linewidth - 14\tabcolsep) * \real{0.1250}}@{}}
\toprule\noalign{}
\endhead
\bottomrule\noalign{}
\endlastfoot
\textbf{指标} & \textbf{POI 总量} & \textbf{物质载体数} & \textbf{非遗数} & \textbf{CMI 均值} & \textbf{OAI 均值} & \textbf{THI 均值} & \textbf{MI 均值} \\
POI 总量 & 1.000 & −0.475 & 0.372 & 0.040 & 0.693 & 0.314 & 0.037 \\
物质载体数 & −0.475 & 1.000 & 0.065 & 0.584 & 0.029 & 0.426 & 0.157 \\
非遗数 & 0.372 & 0.065 & 1.000 & −0.323 & 0.289 & 0.660 & 0.782 \\
CMI 均值 & 0.040 & 0.584 & −0.323 & 1.000 & 0.227 & 0.113 & −0.493 \\
OAI 均值 & 0.693 & 0.029 & 0.289 & 0.227 & 1.000 & 0.702 & 0.242 \\
THI 均值 & 0.314 & 0.426 & 0.660 & 0.113 & 0.702 & 1.000 & 0.758 \\
MI 均值 & 0.037 & 0.157 & 0.782 & −0.493 & 0.242 & 0.758 & 1.000 \\
\end{longtable}
}

镇街级样本只有 7 个，相关系数只能作为提示性结果，不作统计显著性解释。结合空间格局来看，物质载体数与 POI 总量呈中等负相关（r = −0.475），再次说明西南部文化资源和东北部旅游供给之间存在方向性分离。非遗数与 THI 均值、MI 均值呈正相关，可能与西樵、九江等非遗较多镇街同时具有部分景区或文化场馆有关。OAI 均值与 POI 总量、THI 均值呈正相关，提示官方认定和旅游可见度之间存在一定联系，但这一结果不能直接解释为``评级导致热度提升''。

\subsection{政策可干预条件与基础约束条件}\label{ux653fux7b56ux53efux5e72ux9884ux6761ux4ef6ux4e0eux57faux7840ux7ea6ux675fux6761ux4ef6}

根据上述相关分析，本文将潜力释放条件分为政策可干预条件与基础约束条件两类（表 6.6）。前者可以通过规划、运营和公共投入进行改善，后者更适合作为筛选依据和差异化设计背景。

\Needspace{12\baselineskip}
表 6.6　政策可干预条件与基础约束条件的划分

{\def\LTcaptype{none} % do not increment counter
\begin{longtable}[]{@{}
  >{\centering\arraybackslash}p{(\linewidth - 6\tabcolsep) * \real{0.2500}}
  >{\centering\arraybackslash}p{(\linewidth - 6\tabcolsep) * \real{0.2500}}
  >{\centering\arraybackslash}p{(\linewidth - 6\tabcolsep) * \real{0.2500}}
  >{\centering\arraybackslash}p{(\linewidth - 6\tabcolsep) * \real{0.2500}}@{}}
\toprule\noalign{}
\endhead
\bottomrule\noalign{}
\endlastfoot
\textbf{条件类型} & \textbf{典型变量或环节} & \textbf{数据指向} & \textbf{建议着力方向} \\
政策可干预 & 评级申报与名录完善 & OAI 与 POI 总量、THI 均值存在正相关 & 优先支持条件成熟的沉睡潜力载体完善申报和服务标准 \\
政策可干预 & 解说系统与体验内容 & 评分与 THI 关系较强 & 在空心景点区补足地方文化叙事和互动体验 \\
政策可干预 & 游线组织与交通接驳 & POI 数量与 THI 关系较强 & 强化西南片区节点串联和公共交通可达性 \\
政策可干预 & 活态传习与活动运营 & 评论量与 THI 关系较强 & 在历史村落、非遗空间和文化场馆中组织可停留活动 \\
基础约束 & 文化记忆深度 & CMI 相对稳定，难以短期改变 & 作为识别潜力资源的依据 \\
基础约束 & 城镇化和市场规模 & 东北部 POI 供给更集中 & 作为解释旅游热度差异的背景变量 \\
基础约束 & 山水本底与历史空间格局 & 西南部文化资源更集中 & 作为片区化设计和游线组织依据 \\
\end{longtable}
}

这一划分表明，沉睡潜力区的重点不是简单增加景点数量，而是补足可到达、可停留、可解释的条件；空心景点区的重点也不是降低旅游热度，而是把已有客流和场地转化为更有地方文化内容的体验。对于核心耦合区，后续工作不应止于确认其``双高''状态，还应继续检查游客评论、产品内容和解说系统中是否真正出现地方文化识别。

\section{九江片区的专题讨论}\label{ux4e5dux6c5fux7247ux533aux7684ux4e13ux9898ux8ba8ux8bba}

\subsection{反差概况}\label{ux53cdux5deeux6982ux51b5}

九江镇拥有 13 项非遗，并分布有中国历史文化名镇、历史文化名村和多处文物保护单位，是南海区文化资源较为集中的镇街之一。但从旅游供给侧看，九江镇 POI 仅 259 条，占全区 1.92\%，在七个镇街中最低；在当前平台评论口径下，其相关评论量也明显偏低。因此，九江呈现出较典型的``文化线索较密、旅游转化不足''的现象。

\begin{figure}
\centering
\includegraphics[width=5.90556in,height=4.06612in]{ms_thesis/image16.png}
\caption{九江片区错位专题图}
\end{figure}

\subsection{反差的两重来源}\label{ux53cdux5deeux7684ux4e24ux91cdux6765ux6e90}

九江的反差不能只从单一角度解释。一方面，它确实存在旅游转化不足的问题。九江镇位于南海西南部，对外交通和景区化程度相对弱于桂城、大沥等片区，烟桥烟南村等载体的 MI 值明显为负，说明文化记忆和官方认证尚未充分转化为可见旅游热度。

另一方面，数据口径也会放大九江的低估。九江的文化符号中有相当一部分与美食、龙舟和日常节庆相关，如九江双蒸酒、煎堆、鱼花和龙舟宴等；而本文的 POI 体系没有纳入普通餐饮门店，评论数据也主要来自旅游平台，缺少大众点评、美团等餐饮评价数据。对于九江这类餐饮和民俗体验占比较高的片区，当前数据口径会低估其实际生活热度和消费热度。

\subsection{对规划与研究的启示}\label{ux5bf9ux89c4ux5212ux4e0eux7814ux7a76ux7684ux542fux793a}

因此，九江不宜被简单定义为``旅游洼地''，而应被视为``实质性转化不足''和``数据口径低估''共同作用的复合案例。规划上，九江仍需要补足游线组织、展示空间、公共交通和游客服务；研究上，后续也应优先补充餐饮评论、夜间热力和节庆活动数据，以更准确地识别其真实旅游活力。

\section{官方资源扩展与诊断拆分}\label{ux5b98ux65b9ux8d44ux6e90ux6269ux5c55ux4e0eux8bcaux65adux62c6ux5206}

\subsection{官方资源空间化与去重}\label{ux5b98ux65b9ux8d44ux6e90ux7a7aux95f4ux5316ux4e0eux53bbux91cd}

在 165 条核心指数样本之外，本文建立官方资源扩展层，用于检验核心样本是否遗漏了重要文化资源。扩展层共包含 593 条记录，其中不可移动文物 484 处、博物馆 11 家，历史建筑名录、传统村落、历史文化名村、历史文化街区等历史文化资源共 98 条。空间化时，能够定位到具体点位的资源采用点位表达；缺少精确边界的历史文化名镇、名村和传统村落，则使用地名中心或镇街面作为规划分析代理。

\Needspace{12\baselineskip}
表 6.7　官方资源扩展与去重统计

{\def\LTcaptype{none} % do not increment counter
\begin{longtable}[]{@{}
  >{\centering\arraybackslash}p{(\linewidth - 4\tabcolsep) * \real{0.3333}}
  >{\centering\arraybackslash}p{(\linewidth - 4\tabcolsep) * \real{0.3333}}
  >{\centering\arraybackslash}p{(\linewidth - 4\tabcolsep) * \real{0.3333}}@{}}
\toprule\noalign{}
\endhead
\bottomrule\noalign{}
\endlastfoot
\textbf{统计项} & \textbf{数量} & \textbf{说明} \\
官方资源扩展记录 & 593 & 484 处不可移动文物、11 家博物馆、98 条历史文化资源 \\
与既有文化锚点重名 & 96 & 不重复计入扩展锚点 \\
仅匹配到 POI 的资源 & 73 & 可作为后续人工复核对象 \\
纯新增候选资源 & 424 & 与既有锚点和 POI 均未重名 \\
扩展库新增锚点 & 490 & 未与既有锚点重名的官方资源合并后计入 \\
扩展锚点库总量 & 710 & 原 220 条锚点与 490 条新增锚点合计 \\
镇街代表点兜底 & 148 & 空间精度较低，需在后续规划中复核 \\
\end{longtable}
}

\begin{figure}
\centering
\includegraphics[width=5.90556in,height=4.23232in]{ms_thesis/image17.png}
\caption{官方资源扩展空间化图}
\end{figure}

\begin{figure}
\centering
\includegraphics[width=5.90556in,height=4.23232in]{ms_thesis/image18.png}
\caption{500 m 网格官方资源覆盖图}
\end{figure}

图 6.6 和图 6.7 显示，官方资源扩展后，狮山、西樵、大沥、里水等镇街的不可移动文物和历史建筑数量得到补充。其中，西樵镇因中国历史文化名镇采用镇街面表达，官方覆盖网格数较多。这个结果适合作为官方认定范围的上限判断，但不能直接等同于每个网格都具备同等强度的旅游开发条件。

\subsection{典籍---官方---旅游三维诊断}\label{ux5178ux7c4dux5b98ux65b9ux65c5ux6e38ux4e09ux7ef4ux8bcaux65ad}

为解释错位来源，本文按镇街拆分典籍/图谱文化厚度、官方资源覆盖和旅游热度三类指标。该拆分不改变前文错位指数的核心模型，而是用于判断某一片区的短板主要来自文本叙事不足、官方资源覆盖不足，还是旅游转化不足。


\noindent\textbf{表 6.8　镇街典籍---官方---旅游诊断拆分摘要}\par\vspace{0.35em}
\begin{center}
{\small
\setlength{\tabcolsep}{4pt}
\renewcommand{\arraystretch}{1.15}
\begin{tabularx}{\textwidth}{>{\raggedright\arraybackslash}p{0.14\textwidth}
>{\centering\arraybackslash}p{0.15\textwidth}
>{\centering\arraybackslash}p{0.15\textwidth}
>{\centering\arraybackslash}p{0.15\textwidth}
>{\centering\arraybackslash}p{0.16\textwidth}
>{\centering\arraybackslash}X}
\toprule
\textbf{镇街} & \textbf{\makecell{C 典籍/\\图谱均值}} & \textbf{\makecell{O 官方资源\\均值}} & \textbf{\makecell{T 旅游热度\\均值}} & \textbf{\makecell{修正错位\\均值}} & \textbf{\makecell{官方覆盖\\网格}} \\
\midrule
桂城街道 & 2.15 & 2.60 & 22.28 & 19.95 & 40 \\
大沥镇 & 1.73 & 4.01 & 17.18 & 14.54 & 67 \\
里水镇 & 0.37 & 2.09 & 7.59 & 6.53 & 54 \\
狮山镇 & 0.32 & 1.36 & 5.72 & 4.98 & 86 \\
丹灶镇 & 4.36 & 1.84 & 5.07 & 1.71 & 44 \\
西樵镇 & 11.78 & 18.29 & 6.55 & −7.83 & 715 \\
九江镇 & 6.17 & 2.46 & 5.70 & 1.01 & 35 \\
\bottomrule
\end{tabularx}
}
\end{center}

\begin{figure}[H]
\centering
\includegraphics[width=\textwidth,height=0.42\textheight,keepaspectratio]{ms_thesis/image19.png}
\caption{典籍---官方---旅游诊断拆分图}
\end{figure}

表 6.8 和图 6.8 强化了前文判断：桂城、大沥的旅游热度明显高于典籍/图谱和官方资源厚度，属于``旅游先行、文化注入不足''的片区；西樵的典籍/图谱和官方资源均较高，但旅游热度相对偏低，仍是``文化厚、转化不足''的重点片区；九江的典籍/图谱均值高于官方资源覆盖值，说明其问题不仅在旅游转化，也在于部分地方文化线索尚未充分进入可空间化、可体验的资源体系。核心模型用于识别错位强度，三维拆分用于解释错位来源，两者可以互相补充。

从 165 条载体的 CMI、OAI、THI 三个指标看，CMI 与 THI 的 Pearson 相关系数为 0.011，Spearman 相关系数为 0.023，几乎没有相关关系；OAI 与 THI 的相关性也较弱（Pearson = 0.146，Spearman = 0.125）。这说明典籍/图谱中的文化记忆、官方认证和平台旅游热度并不是同一件事，也为后文将``核心耦合区''解释为优先深化潜力而非融合成熟提供了依据。

\section{稳健性与评论语义补充检验}\label{ux7a33ux5065ux6027ux4e0eux8bc4ux8bbaux8bedux4e49ux8865ux5145ux68c0ux9a8c}

\subsection{权重敏感性检验}\label{ux6743ux91cdux654fux611fux6027ux68c0ux9a8c}

前文指数采用 CMI 与 OAI 等权、THI 内部按 POI 数量、平台评分和评论量综合加权的基准口径。为回应权重设定可能带来的主观性问题，本文设置四组替代口径进行敏感性检验：一是提高文化记忆权重，二是提高官方认证权重，三是提高 POI 数量权重，四是提高评论量权重。检验目的不是重新寻找一个``最优权重''，而是观察主要分类结果是否会因合理范围内的权重调整而发生明显改变。


\begin{table}[H]
\centering
\noindent\textbf{表 6.9　权重敏感性检验结果}\par\vspace{0.35em}
{\scriptsize
\setlength{\tabcolsep}{2.2pt}
\renewcommand{\arraystretch}{1.12}
\begin{tabularx}{\textwidth}{@{}>{\raggedright\arraybackslash}p{0.12\textwidth}
>{\raggedright\arraybackslash}X
>{\centering\arraybackslash}p{0.07\textwidth}
>{\centering\arraybackslash}p{0.07\textwidth}
>{\centering\arraybackslash}p{0.07\textwidth}
>{\centering\arraybackslash}p{0.07\textwidth}
>{\centering\arraybackslash}p{0.09\textwidth}@{}}
\toprule
\textbf{口径} & \textbf{权重设置} & \textbf{\makecell{核心\\耦合区}} & \textbf{\makecell{一般\\耦合区}} & \textbf{\makecell{沉睡\\潜力区}} & \textbf{\makecell{空心\\景点区}} & \textbf{\makecell{与基准\\一致比例}} \\
\midrule
基准口径 & CMI/OAI=0.5/0.5；THI=POI 0.4 + 评分 0.2 + 评论 0.4 & 5 & 57 & 54 & 49 & 100.0\% \\
文化侧加权 & CMI/OAI=0.7/0.3 & 5 & 45 & 59 & 56 & 89.7\% \\
官方侧加权 & CMI/OAI=0.3/0.7 & 5 & 70 & 49 & 41 & 90.9\% \\
POI 导向 THI & THI=POI 0.5 + 评分 0.2 + 评论 0.3 & 4 & 59 & 53 & 49 & 97.0\% \\
评论导向 THI & THI=POI 0.3 + 评分 0.2 + 评论 0.5 & 5 & 59 & 53 & 48 & 98.8\% \\
\bottomrule
\end{tabularx}
}
\end{table}

\begin{figure}
\centering
\includegraphics[width=5.90556in,height=\textheight,keepaspectratio]{ms_thesis/fig_6_9_sensitivity_final.png}
\caption{不同权重口径下的错位分类稳定性}
\end{figure}

表 6.9 和图 6.9 显示，四组替代口径下的分类一致比例均在 89.7\% 以上，本文关于``沉睡潜力区''和``空心景点区''并存的总体判断具有一定稳定性。核心耦合区在多数口径下保持 5 条，只有 POI 导向 THI 口径下降为 4 条，说明高值耦合样本数量本身较少，且对旅游热度构成略有敏感。对于在不同权重下发生变化的样本，后续规划更适合结合现场条件、运营基础和游客评论再作复核。

\subsection{评论文化识别度补充检验}\label{ux8bc4ux8bbaux6587ux5316ux8bc6ux522bux5ea6ux62bdux6837ux68c0ux9a8c}

为补充 THI 对游客文化识别解释不足的问题，本文在已匹配的评论数据基础上，对严格匹配的 117 个 POI 进行评论语义抽取，构建以 POI 为中心、以评论实体为外延的轻量评论知识图谱。图谱将评论实体分为体验主题、文化空间、历史记忆、岭南地方性、人物和非遗民俗六类，用于观察游客评论中是否已经出现可识别的文化内容。该检验不改变前文 THI 和 MI 的计算口径，而是作为旅游热度之外的语义校验。

\begin{figure}[H]
\centering
\includegraphics[width=0.92\textwidth]{ms_thesis/fig_6_10_comment_kg_overview.png}
\caption{游客评论文化识别知识图谱}
\end{figure}

图 6.10 中，粉色节点为评论地点，紫色节点为游客评论中识别出的文化实体或体验主题。评论图谱呈现出中心团块与边缘小团块并存的结构，游客评论中的文化识别并非均匀分布，而是集中在少数高识别地点和若干共享主题上。摄影打卡、空间景观等体验主题连接范围较广，也提示旅游热度首先表现为视觉和空间体验，仍需观察其是否转化为人物、非遗、历史记忆和地方性文化识别。


\begin{table}[H]
\centering
\noindent\textbf{表 6.10　游客评论文化识别度补充检验结果}\par\vspace{0.35em}
{\scriptsize
\setlength{\tabcolsep}{3.0pt}
\renewcommand{\arraystretch}{1.15}
\begin{tabularx}{\textwidth}{@{}>{\centering\arraybackslash}p{0.18\textwidth}
>{\centering\arraybackslash}p{0.24\textwidth}
>{\raggedright\arraybackslash}X@{}}
\toprule
\textbf{检验内容} & \textbf{结果} & \textbf{说明} \\
\midrule
图谱规模 & 117 个 POI、327 个评论实体、481 条关系 & 覆盖严格匹配到南海 POI 的游客评论，不将未形成实体的地点强行解释为文化识别。 \\
识别等级 & 高识别 9 个、中识别 30 个、弱识别 15 个、未形成 63 个 & 超过一半 POI 尚未在评论中形成明确文化识别，说明热度不等于文化感知。 \\
实体谱系 & 体验主题 63、文化空间 95、历史记忆 54、岭南地方性 40、人物 38、非遗民俗 37 & 评论中最强的是空间与体验，其次才是人物、非遗、历史和地方性文化。 \\
高识别样本 & 黄飞鸿狮艺武术馆、康有为故居、西樵山、吴家大院、烟桥古村、叶问宗师馆等 & 这些地点更适合作为优先深化对象，而不是简单认定为融合已经成熟。 \\
共享高频主题 & ``摄影打卡''和``空间景观''各连接 32 个地点 & 游客共同感知首先来自视觉和空间体验，文化叙事仍需通过解说、活动和产品进一步转化。 \\
\bottomrule
\end{tabularx}
}
\end{table}

检验结果显示，游客评论中的文化识别并不均衡。黄飞鸿狮艺武术馆、康有为故居、吴家大院、烟桥古村等样本已经形成较清晰的人物、非遗、宗祠、古村或历史记忆识别，可以作为优先深化节点；而大量公园、湿地、城市开放空间和自然景观虽然有游览热度，但评论主要集中在摄影打卡、空间景观、休闲服务等体验主题上，文化内容尚未充分进入游客感知。因此，对于有文化基础但评论识别不足的载体，后续更应改善展示、解说、活动和线上传播条件。
